\documentclass[oneside]{book}

% ------------------------------------------------------------------------------
% Packages
% ------------------------------------------------------------------------------

\usepackage[utf8]{inputenc}
\usepackage[english]{babel}

\usepackage{amsmath, amsthm, amssymb, amsfonts}
\usepackage{mathtools}

\usepackage[dvipsnames]{xcolor}
\usepackage{tcolorbox}
\tcbuselibrary{theorems, skins, breakable}

\usepackage{geometry}
\usepackage{setspace}
\usepackage{hyperref}
\usepackage{microtype}

\usepackage{multicol}
\setlength{\columnsep}{1.2em}

\usepackage{graphicx}

\usepackage{booktabs}
\usepackage{array}
\usepackage{tabularx}
\usepackage{longtable}

\newcolumntype{L}[1]{>{\raggedright\arraybackslash}p{#1}}
\newcolumntype{C}[1]{>{\centering\arraybackslash}p{#1}}
\newcolumntype{R}[1]{>{\raggedleft\arraybackslash}p{#1}}

% ------------------------------------------------------------------------------
% Page Setup
% ------------------------------------------------------------------------------

\geometry{
    top=0.8in,
    bottom=0.8in,
    left=0.7in,
    right=0.7in,
    headheight=12pt,
    headsep=20pt,
    footskip=25pt
}

\setstretch{1.15}

\hypersetup{
    colorlinks=true,
    linkcolor=MidnightBlue,
    urlcolor=RoyalBlue,
    citecolor=ForestGreen
}

% ------------------------------------------------------------------------------
% Variables
% ------------------------------------------------------------------------------

\def\notetitle{Finance \& Economics: What They Expect You to Know}
\def\notesubtitle{Micro, Macro, Accounting, Corporate Finance, Investments, Quant Finance, and Game Theory}
\def\noteauthor{Alan Li}
\def\noteinstitution{New York University}

% ------------------------------------------------------------------------------
% Theorem System and User Commands
% ------------------------------------------------------------------------------

% !TEX root = main.tex

% Theorem System
% The following boxes are provided:
%   Definition:     \defn 
%   Theorem:        \thm 
%   Lemma:          \lem
%   Corollary:      \cor
%   Proposition:    \prop   
%   Claim:          \clm
%   Fact:           \fact
%   Proof:          \pf
%   Example:        \ex
%   Remark:         \rmk (sentence), \rmkb (block)
% Suffix
%   r:              Allow Theorem/Definition to be referenced, e.g. thmr
%   p:              Add a short proof block for Lemma, Corollary, Proposition or Claim, e.g. lemp
%                   For theorems, use \pf for proof blocks

\tcbset{
    kb/common block/.style={
        enhanced,
        breakable,
        frame hidden,
        boxrule=0pt,
        sharp corners,
        boxsep=0pt,
        left=6mm,
        right=5mm,
        top=4mm,
        bottom=4mm,
        before skip=10pt,
        after skip=12pt,
        before upper={\parindent15pt\relax},
    },
    kb/proof block/.style={
        kb/common block,
        colback=black!1!white,
        borderline west={0.9mm}{0pt}{black!55},
    },
    kb/accent proof/.style n args={2}{
        kb/common block,
        colback=#1!4!white,
        borderline west={1mm}{0pt}{#1},
        before upper={
            \parindent15pt\relax
            {\noindent\bfseries #2}\par\smallskip
        },
    },
    kb/example block/.style={
        kb/common block,
        colback=MidnightBlue!3!white,
        borderline west={1.2mm}{0pt}{MidnightBlue!70!black},
        borderline north={0.3mm}{0pt}{MidnightBlue!18!white},
    },
    kb/example title/.style={
        colback=MidnightBlue!75!black,
        colframe=MidnightBlue!75!black,
        boxrule=0pt,
        sharp corners,
        left=3mm,
        right=3mm,
        top=1.2mm,
        bottom=1.2mm,
    },
}

\newcounter{example}[section]
\renewcommand{\theexample}{\thesection.\arabic{example}}

% Definition
\newtcbtheorem[number within=section]{mydefinition}{Definition}
{
    enhanced,
    frame hidden,
    titlerule=0mm,
    toptitle=1mm,
    bottomtitle=1mm,
    fonttitle=\bfseries\large,
    coltitle=black,
    colbacktitle=blue!20!white,
    colback=blue!10!white,
}{defn}

\NewDocumentCommand{\defn}{m+m}{
    \begin{mydefinition}{#1}{}
        #2
    \end{mydefinition}
}

\NewDocumentCommand{\defnr}{mm+m}{
    \begin{mydefinition}{#1}{#2}
        #3
    \end{mydefinition}
}

\NewDocumentEnvironment{definition}{o+b}{
    \IfNoValueTF{#1}{
        \begin{mydefinition}{}{}
    }{
        \begin{mydefinition}{#1}{}
    }
    #2
    \end{mydefinition}
}{}

% Theorem
\newtcbtheorem[use counter from=mydefinition]{mytheorem}{Theorem}
{
    enhanced,
    frame hidden,
    titlerule=0mm,
    toptitle=1mm,
    bottomtitle=1mm,
    fonttitle=\bfseries\large,
    coltitle=black,
    colbacktitle=cyan!20!white,
    colback=cyan!10!white,
}{thm}

\NewDocumentCommand{\thm}{m+m}{
    \begin{mytheorem}{#1}{}
        #2
    \end{mytheorem}
}

\NewDocumentCommand{\thmr}{mm+m}{
    \begin{mytheorem}{#1}{#2}
        #3
    \end{mytheorem}
}

\NewDocumentEnvironment{theorem}{o+b}{
    \IfNoValueTF{#1}{
        \begin{mytheorem}{}{}
    }{
        \begin{mytheorem}{#1}{}
    }
    #2
    \end{mytheorem}
}{}

% Lemma
\newtcbtheorem[use counter from=mydefinition]{mylemma}{Lemma}
{
    enhanced,
    frame hidden,
    titlerule=0mm,
    toptitle=1mm,
    bottomtitle=1mm,
    fonttitle=\bfseries\large,
    coltitle=black,
    colbacktitle=violet!20!white,
    colback=violet!10!white,
}{lem}

\NewDocumentCommand{\lem}{m+m}{
    \begin{mylemma}{#1}{}
        #2
    \end{mylemma}
}

\newenvironment{lempf}{
    \begin{tcolorbox}[kb/accent proof={violet!65!black}{Proof for Lemma.}]
}{
    \hfill\textcolor{violet!65!black}{$\blacksquare$}
    \end{tcolorbox}
}

\NewDocumentCommand{\lemp}{m+m+m}{
    \begin{mylemma}{#1}{}
        #2
    \end{mylemma}

    \begin{lempf}
        #3
    \end{lempf}
}

% Corollary
\newtcbtheorem[use counter from=mydefinition]{mycorollary}{Corollary}
{
    enhanced,
    frame hidden,
    titlerule=0mm,
    toptitle=1mm,
    bottomtitle=1mm,
    fonttitle=\bfseries\large,
    coltitle=black,
    colbacktitle=orange!20!white,
    colback=orange!10!white,
}{cor}

\NewDocumentCommand{\cor}{+m}{
    \begin{mycorollary}{}{}
        #1
    \end{mycorollary}
}

\newenvironment{corpf}{
    \begin{tcolorbox}[kb/accent proof={orange!80!black}{Proof for Corollary.}]
}{
    \hfill\textcolor{orange!80!black}{$\blacksquare$}
    \end{tcolorbox}
}

\NewDocumentCommand{\corp}{m+m+m}{
    \begin{mycorollary}{}{}
        #1
    \end{mycorollary}

    \begin{corpf}
        #2
    \end{corpf}
}

% Proposition
\newtcbtheorem[use counter from=mydefinition]{myproposition}{Proposition}
{
    enhanced,
    frame hidden,
    titlerule=0mm,
    toptitle=1mm,
    bottomtitle=1mm,
    fonttitle=\bfseries\large,
    coltitle=black,
    colbacktitle=yellow!30!white,
    colback=yellow!20!white,
}{prop}

\NewDocumentCommand{\prop}{+m}{
    \begin{myproposition}{}{}
        #1
    \end{myproposition}
}

\newenvironment{proppf}{
    \begin{tcolorbox}[kb/accent proof={Goldenrod!85!black}{Proof for Proposition.}]
}{
    \hfill\textcolor{Goldenrod!85!black}{$\blacksquare$}
    \end{tcolorbox}
}

\NewDocumentCommand{\propp}{+m+m}{
    \begin{myproposition}{}{}
        #1
    \end{myproposition}

    \begin{proppf}
        #2
    \end{proppf}
}

% Claim
\newtcbtheorem[use counter from=mydefinition]{myclaim}{Claim}
{
    enhanced,
    frame hidden,
    titlerule=0mm,
    toptitle=1mm,
    bottomtitle=1mm,
    fonttitle=\bfseries\large,
    coltitle=black,
    colbacktitle=pink!30!white,
    colback=pink!20!white,
}{clm}


\NewDocumentCommand{\clm}{m+m}{
    \begin{myclaim*}{#1}{}
        #2
    \end{myclaim*}
}

\newenvironment{clmpf}{
    \begin{tcolorbox}[kb/accent proof={Magenta!70!black}{Proof for Claim.}]
}{
    \hfill\textcolor{Magenta!70!black}{$\blacksquare$}
    \end{tcolorbox}
}

\NewDocumentCommand{\clmp}{m+m+m}{
    \begin{myclaim*}{#1}{}
        #2
    \end{myclaim*}

    \begin{clmpf}
        #3
    \end{clmpf}
}

% Fact
\newtcbtheorem[use counter from=mydefinition]{myfact}{Fact}
{
    enhanced,
    frame hidden,
    titlerule=0mm,
    toptitle=1mm,
    bottomtitle=1mm,
    fonttitle=\bfseries\large,
    coltitle=black,
    colbacktitle=purple!20!white,
    colback=purple!10!white,
}{fact}

\NewDocumentCommand{\fact}{+m}{
    \begin{myfact}{}{}
        #1
    \end{myfact}
}


% Proof
\RenewDocumentEnvironment{proof}{O{\proofname}}{
    \par
    \pushQED{\qed}
    \begin{tcolorbox}[kb/proof block]
        {\noindent\bfseries #1}\par\smallskip
        \ignorespaces
}{
    \popQED
    \end{tcolorbox}
    \par
}

\let\proofenvironment\proof
\let\endproofenvironment\endproof

\makeatletter
\renewcommand{\proof}{
    \@ifnextchar\bgroup{\proofcommand}{\proofenvironment}
}
\makeatother

\NewDocumentCommand{\proofcommand}{+m}{
    \proofenvironment
        #1
    \endproofenvironment
}

\NewDocumentCommand{\pf}{o+m}{
    \IfNoValueTF{#1}{
        \begin{proof}
            #2
        \end{proof}
    }{
        \begin{proof}[#1]
            #2
        \end{proof}
    }
}

% Example
\NewDocumentEnvironment{example}{o+b}{
    \refstepcounter{example}
    \IfNoValueTF{#1}{
        \def\exampletitle{Example \theexample}
    }{
        \def\exampletitle{Example \theexample \textbar\ #1}
    }
    \begin{tcolorbox}[
        kb/example block,
        title={\exampletitle},
        fonttitle=\bfseries\small,
        coltitle=white,
        boxed title style={kb/example title},
        attach boxed title to top left={xshift=4mm,yshift*=-\tcboxedtitleheight/2},
        top=7mm,
    ]
        #2
    \end{tcolorbox}
}{}

\NewDocumentCommand{\ex}{+m}{
    \begin{example}
        #1
    \end{example}
}


% Remark
\NewDocumentCommand{\rmk}{+m}{
    {\it \color{blue!50!white}#1}
}

\newenvironment{remark}{
    \par
    \vspace{5pt}
    \begin{minipage}{\textwidth}
        {\par\noindent{\textbf{Remark.}}}
        \tcolorbox[blanker,breakable,left=5mm,
        before skip=10pt,after skip=10pt,
        borderline west={1mm}{0pt}{cyan!10!white}]
}{
        \endtcolorbox
    \end{minipage}
    \vspace{5pt}
}

\NewDocumentCommand{\rmkb}{+m}{
    \begin{remark}
        #1
    \end{remark}
}

% !TEX root = main.tex

\newcommand{\R}{\mathbb{R}}
\newcommand{\N}{\mathbb{N}}
\newcommand{\vect}[1]{\mathbf{#1}}
\newcommand{\norm}[1]{\left\lVert #1 \right\rVert}
\newcommand{\dd}{\mathop{}\!\mathrm{d}}
\newcommand{\grad}{\nabla}
\newcommand{\diver}{\operatorname{div}}
\newcommand{\pd}[2]{\frac{\partial #1}{\partial #2}}
\newcommand{\lcm}{\operatorname{lcm}}


% ------------------------------------------------------------------------------
% Document
% ------------------------------------------------------------------------------

\begin{document}

% ------------------------------------------------------------------------------
% Title Page
% ------------------------------------------------------------------------------

\begin{titlepage}
    \centering

    \vspace*{2.2cm}

    {\Huge\bfseries \notetitle \par}

    \vspace{0.9cm}

    \resizebox{0.88\textwidth}{!}{%
        \itshape \notesubtitle
    }

    \vspace{2.4cm}

    \vspace{1.2cm}

    {\large \noteauthor \par}

    \vspace{0.25cm}

    {\large \noteinstitution \par}

    \vfill
\end{titlepage}

% ------------------------------------------------------------------------------
% Front Matter
% ------------------------------------------------------------------------------

\frontmatter
\setcounter{tocdepth}{1}
\tableofcontents
\newpage

% ------------------------------------------------------------------------------
% Main Matter
% ------------------------------------------------------------------------------

\mainmatter

% ==============================================================================
\part{Microeconomics}
% ==============================================================================

% ------------------------------------------------------------------------------
\chapter{Consumer Theory}
% ------------------------------------------------------------------------------

\section{Preferences and Utility}

\defn{Preference Relation}{
    A preference relation $\succeq$ on a consumption set $X \subseteq \mathbb{R}^n_+$
    satisfies:
    \begin{itemize}
        \item \textbf{Completeness}: for all $x, y \in X$, either $x \succeq y$ or
              $y \succeq x$ (or both).
        \item \textbf{Transitivity}: if $x \succeq y$ and $y \succeq z$, then $x \succeq z$.
    \end{itemize}
    We write $x \succ y$ (strict preference) if $x \succeq y$ but not $y \succeq x$, and
    $x \sim y$ (indifference) if both $x \succeq y$ and $y \succeq x$.
}

\defn{Utility Function}{
    A function $u: X \to \mathbb{R}$ \textbf{represents} the preference relation
    $\succeq$ if for all $x, y \in X$:
    \[
        x \succeq y \iff u(x) \geq u(y).
    \]
    Utility is \textbf{ordinal}: any strictly increasing transformation $v = f(u)$
    with $f' > 0$ represents the same preferences.
}

\rmkb{Common utility functions for two goods $(x_1, x_2)$:
\begin{center}
\begin{tabularx}{\linewidth}{lllX}
\toprule
\textbf{Type} & \textbf{$u(x_1,x_2)$} & \textbf{MRS} & \textbf{Substitutability} \\
\midrule
Cobb-Douglas & $x_1^\alpha x_2^\beta$ & $\frac{\alpha x_2}{\beta x_1}$ & Diminishing \\
Perfect substitutes & $ax_1 + bx_2$ & $a/b$ (constant) & Perfect \\
Perfect complements & $\min\{ax_1, bx_2\}$ & $0$ or $\infty$ & Zero \\
Quasilinear & $v(x_1) + x_2$ & $v'(x_1)$ & No income effect on $x_1$ \\
CES & $(x_1^\rho + x_2^\rho)^{1/\rho}$ & $(x_2/x_1)^{1-\rho}$ & $\sigma = 1/(1-\rho)$ \\
\bottomrule
\end{tabularx}
\end{center}}

\defn{Marginal Rate of Substitution}{
    The MRS at a bundle $(x_1, x_2)$ is the rate at which the consumer is willing
    to trade $x_2$ for $x_1$ along an indifference curve:
    \[
        \text{MRS}_{12} = -\frac{dx_2}{dx_1}\bigg|_{u=\bar{u}}
        = \frac{\partial u/\partial x_1}{\partial u/\partial x_2}
        = \frac{MU_1}{MU_2}.
    \]
    Convex preferences imply a \textbf{diminishing} MRS.
}

% --------------------------------------------------------------------------
\section{Budget Constraint and Consumer Optimum}
% --------------------------------------------------------------------------

\defn{Budget Constraint}{
    With prices $(p_1, p_2)$ and income $m$, the budget set is
    \[
        \mathcal{B}(p, m) = \{(x_1, x_2) \in \mathbb{R}^2_+ : p_1 x_1 + p_2 x_2 \leq m\}.
    \]
    The budget line $p_1 x_1 + p_2 x_2 = m$ has slope $-p_1/p_2$ and intercepts
    $m/p_1$ and $m/p_2$.
}

\thm{Consumer Optimum (Interior Solution)}{
    If preferences are strictly convex and differentiable, the optimal bundle
    $(x_1^*, x_2^*)$ satisfies the budget constraint with equality and
    \[
        \text{MRS}_{12} = \frac{p_1}{p_2}
        \qquad \Longleftrightarrow \qquad
        \frac{MU_1}{MU_2} = \frac{p_1}{p_2}
        \qquad \Longleftrightarrow \qquad
        \frac{MU_1}{p_1} = \frac{MU_2}{p_2}.
    \]
    The last form says that at the optimum, the marginal utility per dollar spent
    is equalised across all goods (\textbf{equal bang per buck}).
}

\pf{
    The Lagrangian is $\mathcal{L} = u(x_1, x_2) - \lambda(p_1 x_1 + p_2 x_2 - m)$.
    First-order conditions: $\partial u/\partial x_i = \lambda p_i$ for $i = 1, 2$.
    Dividing: $(\partial u/\partial x_1)/(\partial u/\partial x_2) = p_1/p_2$.
}

% --------------------------------------------------------------------------
\section{Demand Functions}
% --------------------------------------------------------------------------

\defn{Marshallian (Ordinary) Demand}{
    The solution to $\max_{x} u(x)$ s.t.\ $p \cdot x = m$ gives the Marshallian
    demand functions $x_i^*(p_1, p_2, m)$.  For Cobb-Douglas $u = x_1^\alpha x_2^\beta$:
    \[
        x_1^* = \frac{\alpha}{\alpha+\beta}\cdot\frac{m}{p_1}, \qquad
        x_2^* = \frac{\beta}{\alpha+\beta}\cdot\frac{m}{p_2}.
    \]
    The expenditure share on good $i$ is constant: $p_i x_i^*/m = \alpha_i/(\alpha+\beta)$.
}

\defn{Hicksian (Compensated) Demand}{
    The solution to $\min_{x} p \cdot x$ s.t.\ $u(x) \geq \bar{u}$ gives the
    Hicksian demand $h_i(p, \bar{u})$.  By \textbf{Shephard's Lemma}:
    \[
        h_i(p, \bar{u}) = \frac{\partial e(p, \bar{u})}{\partial p_i},
    \]
    where $e(p, \bar{u})$ is the \textbf{expenditure function} (minimum cost to
    achieve utility $\bar{u}$ at prices $p$).
}

\thm{Slutsky Equation}{
    The total price effect decomposes into a substitution effect and an income effect:
    \[
        \underbrace{\frac{\partial x_i^*}{\partial p_j}}_{\text{total}}
        = \underbrace{\frac{\partial h_i}{\partial p_j}}_{\substack{\text{substitution}\\\text{(always }\leq 0\text{ if }i=j\text{)}}}
        - \underbrace{x_j^* \frac{\partial x_i^*}{\partial m}}_{\text{income}}.
    \]
    For a \textbf{normal good}, both effects reinforce (demand falls as $p_i$ rises).
    For a \textbf{Giffen good}, the income effect dominates and demand rises with price.
}

\rmkb{The Slutsky matrix $S_{ij} = \partial h_i/\partial p_j$ is symmetric and
negative semi-definite.  Own-price substitution effects are always non-positive:
$S_{ii} \leq 0$.}

% --------------------------------------------------------------------------
\section{Revealed Preference and Welfare}
% --------------------------------------------------------------------------

\defn{Indirect Utility and Roy's Identity}{
    The \textbf{indirect utility function} $v(p, m) = u(x^*(p, m))$ gives maximum
    utility as a function of prices and income.  \textbf{Roy's Identity} recovers
    Marshallian demand from $v$:
    \[
        x_i^*(p, m) = -\frac{\partial v/\partial p_i}{\partial v/\partial m}.
    \]
}

\defn{Equivalent and Compensating Variation}{
    For a price change from $p^0$ to $p^1$:
    \[
        CV = e(p^1, u^0) - e(p^0, u^0) = m - e(p^1, u^0),
    \]
    \[
        EV = e(p^1, u^1) - e(p^0, u^1) = e(p^0, u^1) - m.
    \]
    Both measure welfare change in monetary terms; they coincide for quasilinear
    preferences (where consumer surplus is an exact welfare measure).
}

% ==============================================================================
\chapter{Producer Theory}
% ==============================================================================

\section{Production Technology}

\defn{Production Function}{
    $f(K, L)$ maps capital $K$ and labour $L$ to maximum output $q$.
    The \textbf{marginal product} of input $i$ is $MP_i = \partial f/\partial x_i > 0$.
    \textbf{Diminishing marginal product}: $\partial^2 f/\partial x_i^2 < 0$.
}

\defn{Returns to Scale}{
    For scalar $\lambda > 1$:
    \begin{itemize}
        \item \textbf{CRS} (Constant): $f(\lambda K, \lambda L) = \lambda f(K,L)$.
        \item \textbf{IRS} (Increasing): $f(\lambda K, \lambda L) > \lambda f(K,L)$.
        \item \textbf{DRS} (Decreasing): $f(\lambda K, \lambda L) < \lambda f(K,L)$.
    \end{itemize}
    Cobb-Douglas $f = K^\alpha L^\beta$: CRS iff $\alpha + \beta = 1$.
}

\defn{Marginal Rate of Technical Substitution}{
    The MRTS is the slope of an isoquant:
    \[
        \text{MRTS}_{KL} = -\frac{dK}{dL}\bigg|_{q=\bar{q}}
        = \frac{MP_L}{MP_K}.
    \]
    Diminishing MRTS corresponds to convex isoquants (standard assumption).
}

% --------------------------------------------------------------------------
\section{Cost Minimisation}
% --------------------------------------------------------------------------

\thm{Cost-Minimising Input Choice}{
    The firm solves $\min_{K,L} wL + rK$ s.t.\ $f(K,L) = q$.
    Interior solution satisfies:
    \[
        \frac{MP_L}{w} = \frac{MP_K}{r}
        \qquad \Longleftrightarrow \qquad
        \text{MRTS}_{KL} = \frac{w}{r}.
    \]
}

\defn{Cost Function}{
    The \textbf{total cost function} $c(q, w, r)$ is the minimum expenditure to
    produce $q$ units.  By \textbf{Shephard's Lemma}:
    \[
        \frac{\partial c}{\partial w} = L^*(q,w,r), \qquad
        \frac{\partial c}{\partial r} = K^*(q,w,r).
    \]
    Cost components:
    \[
        TC = FC + VC(q), \quad
        AC(q) = \frac{TC}{q}, \quad
        MC(q) = \frac{dTC}{dq}.
    \]
    The $MC$ curve intersects the $AC$ curve at the minimum of $AC$.
}

\rmkb{For Cobb-Douglas $f = K^\alpha L^\beta$ with $\alpha + \beta = 1$:
$c(q,w,r) = q \cdot \text{const} \cdot w^\beta r^\alpha$, so $MC = AC = $ constant
(CRS $\Rightarrow$ flat LRAC).}

% --------------------------------------------------------------------------
\section{Profit Maximisation}
% --------------------------------------------------------------------------

\defn{Profit Function}{
    $\pi = TR - TC = pq - c(q)$.  Under perfect competition, $p$ is taken as given.
    FOC: $p = MC(q^*)$.  SOC: $MC'(q^*) > 0$ (rising MC).
    The firm shuts down if $p < \min AVC$ (short run) or $p < \min AC$ (long run).
}

\thm{Hotelling's Lemma}{
    If $\pi(p, w, r) = \max_{q,K,L}[pf(K,L) - wL - rK]$, then:
    \[
        \frac{\partial \pi}{\partial p} = q^*, \qquad
        \frac{\partial \pi}{\partial w} = -L^*, \qquad
        \frac{\partial \pi}{\partial r} = -K^*.
    \]
}

% ==============================================================================
\chapter{Market Structures}
% ==============================================================================

\section{Perfect Competition}

\defn{Competitive Equilibrium}{
    Price $p^*$ and quantity $Q^*$ such that $Q^D(p^*) = Q^S(p^*)$.
    Each firm takes $p^*$ as given and produces where $p^* = MC(q_i)$.
    In long-run equilibrium: $p^* = MC = \min AC$ (zero economic profit).
}

\defn{Consumer and Producer Surplus}{
    \[
        CS = \int_0^{Q^*} [P^D(Q) - p^*]\,dQ, \qquad
        PS = \int_0^{Q^*} [p^* - P^S(Q)]\,dQ.
    \]
    \textbf{Total Surplus} $= CS + PS$ is maximised at competitive equilibrium
    (\textbf{First Welfare Theorem}).
}

% --------------------------------------------------------------------------
\section{Monopoly}
% --------------------------------------------------------------------------

\defn{Monopoly Optimum}{
    The monopolist maximises $\pi = TR - TC = p(Q) \cdot Q - c(Q)$.
    FOC: $MR(Q^*) = MC(Q^*)$, where $MR = p + p'(Q)Q = p(1 - 1/|\varepsilon|)$.
    The \textbf{Lerner Index} measures market power:
    \[
        L = \frac{p - MC}{p} = \frac{1}{|\varepsilon_D|}.
    \]
    Deadweight loss (DWL) equals the welfare lost relative to the competitive outcome.
}

\rmkb{\textbf{Price discrimination}: (1st degree) personalized prices extracting
all CS; (2nd degree) nonlinear pricing/quantity discounts; (3rd degree) separate
markets segmented by $\varepsilon$.  Under 3rd degree: $MR_1 = MR_2 = MC$.}

% --------------------------------------------------------------------------
\section{Oligopoly}
% --------------------------------------------------------------------------

\defn{Cournot Competition}{
    $n$ firms simultaneously choose quantities $q_i$.  Firm $i$ maximises:
    \[
        \pi_i = p(Q) q_i - c_i(q_i), \quad Q = \sum_j q_j.
    \]
    FOC (best response): $p(Q) + p'(Q) q_i = MC_i$.
    Symmetric equilibrium with linear demand $p = a - bQ$ and constant $MC$:
    \[
        q^* = \frac{a - MC}{b(n+1)}, \quad p^* = \frac{a + n \cdot MC}{n+1}.
    \]
    As $n \to \infty$: $p^* \to MC$ (competitive outcome).
}

\defn{Bertrand Competition}{
    Firms simultaneously set prices.  With homogeneous goods and constant identical
    $MC$: unique Nash equilibrium is $p_1 = p_2 = MC$ (Bertrand paradox).
    With differentiated products or capacity constraints, equilibrium prices exceed $MC$.
}

\defn{Stackelberg Competition}{
    Leader (firm 1) commits to quantity $q_1$ first; follower (firm 2) best-responds.
    Backward induction: leader internalises the follower's reaction function
    $q_2^{BR}(q_1)$.  Leader has first-mover advantage: higher profit than Cournot.
}

% ==============================================================================
\chapter{Welfare Economics and Market Failures}
% ==============================================================================

\section{Welfare Analysis}

\thm{First Welfare Theorem}{
    If all markets are perfectly competitive and there are no externalities or
    public goods, then every competitive equilibrium is \textbf{Pareto efficient}:
    no reallocation can make someone better off without making someone else worse off.
}

\thm{Second Welfare Theorem}{
    Under convexity assumptions, every Pareto-efficient allocation can be supported
    as a competitive equilibrium with appropriate lump-sum transfers of endowments.
}

% --------------------------------------------------------------------------
\section{Externalities and Public Goods}
% --------------------------------------------------------------------------

\defn{Externality}{
    A cost or benefit imposed on a third party not involved in the transaction.
    \begin{itemize}
        \item \textbf{Negative externality}: $MSC > MPC$ (e.g.\ pollution).
              Optimal Pigouvian tax $t^* = MSC - MPC$ at efficient output.
        \item \textbf{Positive externality}: $MSB > MPB$ (e.g.\ vaccination).
              Optimal Pigouvian subsidy $s^* = MSB - MPB$.
    \end{itemize}
}

\thm{Coase Theorem}{
    If property rights are well-defined and bargaining is costless, private
    negotiation leads to the efficient outcome regardless of the initial assignment
    of rights.  In practice, transaction costs, incomplete information, and
    hold-up problems limit its applicability.
}

\defn{Public Good}{
    A good that is \textbf{non-excludable} (cannot prevent consumption) and
    \textbf{non-rival} (one person's consumption does not reduce another's).
    Samuelson rule for efficient provision: $\sum_i MRS_i = MRT$ (sum of marginal
    willingness to pay equals marginal cost of provision).
    Market under-provides due to the \textbf{free-rider problem}.
}

% ==============================================================================
\chapter{Information Economics}
% ==============================================================================

\section{Adverse Selection}

\defn{Adverse Selection}{
    Arises when one party has \textbf{hidden information} (pre-contractual), causing
    a market to attract a disproportionately unfavourable pool.
    Classic example (Akerlof): in the used-car market with quality $\theta \sim U[0,1]$,
    sellers know $\theta$, buyers do not.  At any pooling price $p$, only sellers
    with $\theta \leq p/v_s$ sell; buyers anticipate average quality $p/(2v_s)$,
    driving the market toward collapse.
}

\defn{Signalling Equilibrium}{
    The high-type agent undertakes a costly signal $s$ to credibly distinguish
    itself.  Spence education model: worker of type $\theta \in \{H, L\}$ chooses
    education $e$.  Separating equilibrium requires the single-crossing (Spence-Mirrlees)
    condition: $\partial^2 u / (\partial e\, \partial \theta) > 0$, ensuring the
    incentive-compatibility constraints are not both binding simultaneously.
    \[
        IC_H: u(w_H, e_H; H) \geq u(w_L, e_L; H), \qquad
        IC_L: u(w_L, e_L; L) \geq u(w_H, e_H; L).
    \]
}

% --------------------------------------------------------------------------
\section{Moral Hazard}
% --------------------------------------------------------------------------

\defn{Moral Hazard}{
    Arises when one party's \textbf{hidden action} (post-contractual) cannot be
    observed or verified.  Principal--agent problem: principal maximises
    $E[q(e)] - w(q)$ subject to the agent's
    \[
        IR: \quad E[u(w(q))] - c(e) \geq \bar{u} \quad \text{(participation)},
    \]
    \[
        IC: \quad e^* = \arg\max_{e}\, E[u(w(q))] - c(e) \quad \text{(incentive compatibility)}.
    \]
    First-best requires both risk sharing and incentive provision; with risk-averse
    agents these goals conflict (\textbf{second-best} involves incomplete insurance).
}

\rmkb{Key results: (1) If agent is risk-neutral, sell the firm (franchise contract)
achieves first-best. (2) Optimal second-best contract balances insurance against
incentives; bonus/commission schemes approximate it. (3) Relative performance
evaluation (RPE) removes common shocks from incentive pay.}

% ==============================================================================
\chapter{Choice Under Uncertainty}
% ==============================================================================

\section{Expected Utility Theory}

\defn{Von Neumann--Morgenstern Expected Utility}{
    A lottery $L = (p_1, x_1; \ldots; p_n, x_n)$ assigns probability $p_i$ to outcome $x_i$.
    Under the VNM axioms (completeness, transitivity, continuity, independence), there
    exists a Bernoulli utility function $u: \mathbb{R} \to \mathbb{R}$ such that:
    \[
        U(L) = \sum_{i=1}^n p_i\, u(x_i) = E[u(x)].
    \]
    $u$ is unique up to positive affine transformations.
}

\defn{Risk Aversion}{
    An agent is \textbf{risk-averse} if $u$ is strictly concave: $u'' < 0$.
    \begin{itemize}
        \item \textbf{Certainty equivalent} $CE$: $u(CE) = E[u(x)]$, so $CE \leq E[x]$.
        \item \textbf{Risk premium} $\pi_R = E[x] - CE \geq 0$.
        \item \textbf{Arrow-Pratt absolute risk aversion}: $A(x) = -u''(x)/u'(x)$.
        \item \textbf{Relative risk aversion}: $R(x) = -xu''(x)/u'(x)$.
    \end{itemize}
    For CARA: $u(x) = -e^{-\alpha x}$, so $A(x) = \alpha$ (constant).
    For CRRA: $u(x) = x^{1-\rho}/(1-\rho)$, so $R(x) = \rho$ (constant).
}

\thm{Jensen's Inequality}{
    For a concave function $u$ and random variable $x$:
    $u(E[x]) \geq E[u(x)]$.
    Equality holds iff $x$ is degenerate (no risk).
}

% ==============================================================================
\part{Macroeconomics}
% ==============================================================================

% ------------------------------------------------------------------------------
\chapter{National Income and Output}
% ------------------------------------------------------------------------------

\section{GDP and National Accounts}

\defn{Gross Domestic Product}{
    GDP measures the total market value of all \textbf{final goods and services}
    produced within a country's borders in a given period.  Three equivalent
    approaches:
    \[
        GDP = C + I + G + NX \quad \text{(expenditure approach)},
    \]
    where $C$ = private consumption, $I$ = gross investment,
    $G$ = government spending, $NX = X - M$ = net exports.
}

\defn{Real vs.\ Nominal GDP}{
    \textbf{Nominal GDP} values output at current prices.
    \textbf{Real GDP} values output at base-year prices, removing the effect of
    price changes.
    \[
        \text{GDP Deflator} = \frac{\text{Nominal GDP}}{\text{Real GDP}} \times 100.
    \]
    Growth in real GDP $\approx$ growth in nominal GDP $-$ inflation.
}

\rmkb{GDP captures production but misses distribution, non-market activity,
informal economy, and welfare.  GNP adds net factor income from abroad.}

% --------------------------------------------------------------------------
\section{Inflation and Price Levels}
% --------------------------------------------------------------------------

\defn{CPI and Inflation}{
    The Consumer Price Index tracks the cost of a fixed basket of goods:
    \[
        \text{CPI}_t = \frac{\sum p_t^i q_0^i}{\sum p_0^i q_0^i} \times 100.
    \]
    \textbf{Inflation rate}: $\pi_t = (CPI_t - CPI_{t-1})/CPI_{t-1}$.
    \textbf{Fisher equation} linking nominal rate $i$, real rate $r$, and inflation $\pi$:
    \[
        (1 + i) = (1 + r)(1 + \pi) \approx r + \pi \quad \text{(for small rates)}.
    \]
}

% ------------------------------------------------------------------------------
\chapter{Money and Monetary Policy}
% ------------------------------------------------------------------------------

\section{Money Supply and Demand}

\defn{Money and its Functions}{
    Money serves as (1) medium of exchange, (2) unit of account, (3) store of value.
    \textbf{M1} = currency + demand deposits.
    \textbf{M2} = M1 + savings deposits + small time deposits.
    \textbf{Money multiplier}: $m = 1/\text{reserve ratio}$ (simple model);
    in practice, $m = (1 + c)/(c + \theta)$ where $c$ = currency-to-deposit ratio,
    $\theta$ = reserve ratio.
}

\section{Monetary Policy and the Interest Rate}

\rmkb{Central banks (e.g., the Fed) target the short-term nominal interest rate
(federal funds rate) using open market operations (buying/selling government bonds).
Expansionary policy: buy bonds $\Rightarrow$ money supply $\uparrow$ $\Rightarrow$
$i \downarrow$ $\Rightarrow$ investment $\uparrow$ $\Rightarrow$ GDP $\uparrow$.
Key constraint: zero lower bound (ZLB); unconventional tools include QE (asset
purchases) and forward guidance.}

% ------------------------------------------------------------------------------
\chapter{Fiscal Policy}
% ------------------------------------------------------------------------------

\section{Government Budget and Multiplier Effects}

\defn{Fiscal Multiplier}{
    An increase in government spending $\Delta G$ raises GDP by more than $\Delta G$:
    \[
        \Delta Y = \frac{1}{1-MPC} \cdot \Delta G = \frac{1}{MPS} \cdot \Delta G,
    \]
    where $MPC$ is the marginal propensity to consume and $MPS = 1 - MPC$.
    Tax multiplier: $\Delta Y = -\frac{MPC}{1-MPC} \cdot \Delta T$ (smaller in magnitude;
    negative sign because tax cut raises disposable income indirectly).
    \textbf{Balanced budget multiplier} $= 1$ (equal $\Delta G$ and $\Delta T$).
}

\rmkb{Crowding out: expansionary fiscal policy may raise interest rates, reducing
private investment.  In a liquidity trap ($i$ at ZLB), crowding out is absent
and the multiplier is larger.  Long-run: deficit spending accumulates debt
($D/GDP$ ratio); sustainability requires $r < g$ (interest rate below growth rate).}

% ------------------------------------------------------------------------------
\chapter{International Economics}
% ------------------------------------------------------------------------------

\section{Exchange Rates and the Balance of Payments}

\defn{Nominal and Real Exchange Rate}{
    The \textbf{nominal exchange rate} $e$ is units of domestic currency per unit
    of foreign currency.  The \textbf{real exchange rate}:
    \[
        q = e \cdot \frac{P^*}{P},
    \]
    where $P^*$ is the foreign price level and $P$ the domestic price level.
    A rise in $q$ is a \textbf{real depreciation} (domestic goods become cheaper).
    \textbf{Purchasing Power Parity (PPP)}: $q = 1$ in the long run, implying
    $e = P/P^*$.
}

\defn{Balance of Payments}{
    \[
        \underbrace{CA}_{\text{Current Account}} + \underbrace{FA}_{\text{Financial Account}} = 0.
    \]
    $CA = NX + \text{net factor income} + \text{net transfers}$.
    A current account \textbf{deficit} means the country imports more than it exports
    and is a net borrower from the rest of the world.
}

\rmkb{Under fixed exchange rates, the central bank must intervene (buy/sell FX
reserves) to maintain the peg; monetary policy autonomy is lost
(\textbf{impossible trinity}: cannot simultaneously have fixed rate, free capital
flows, and independent monetary policy).}

% ==============================================================================
\part{Financial Accounting}
% ==============================================================================

% ------------------------------------------------------------------------------
\chapter{Reading Financial Statements}
% ------------------------------------------------------------------------------

\section{The Three Core Statements}

\defn{Income Statement}{
    Reports \textbf{revenues} and \textbf{expenses} over a period, resulting in
    net income (profit or loss):
    \[
        \text{Net Income} = \text{Revenue} - \text{COGS} - \text{Operating Expenses}
        - \text{Interest} - \text{Taxes}.
    \]
    Key line items: Revenue, Cost of Goods Sold (COGS), \textbf{Gross Profit}
    $=$ Revenue $-$ COGS, \textbf{EBITDA}, \textbf{EBIT} (Operating Income),
    \textbf{EBT}, \textbf{Net Income}.
}

\defn{Balance Sheet (Statement of Financial Position)}{
    A snapshot of what the firm \textbf{owns} (assets), \textbf{owes} (liabilities),
    and the residual claim of shareholders (equity) at a point in time:
    \[
        \text{Assets} = \text{Liabilities} + \text{Shareholders' Equity}.
    \]
    \textbf{Current assets} (cash, AR, inventory) vs.\ \textbf{non-current assets}
    (PP\&E, intangibles, goodwill).
    \textbf{Current liabilities} (AP, short-term debt) vs.\ \textbf{long-term debt}.
}

\defn{Statement of Cash Flows}{
    Reconciles net income to actual cash generated:
    \[
        \text{Cash Flow} = \underbrace{CFO}_{\text{Operations}} + \underbrace{CFI}_{\text{Investing}} + \underbrace{CFF}_{\text{Financing}}.
    \]
    \textbf{CFO}: net income $+$ D\&A $-$ $\Delta$Working Capital (indirect method).
    \textbf{CFI}: capex, acquisitions, asset disposals.
    \textbf{CFF}: debt issuance/repayment, equity issuance, dividends.
    \textbf{Free Cash Flow}: $FCF = CFO - \text{Capex}$.
}

% --------------------------------------------------------------------------
\section{Key Ratios for Analysts}
% --------------------------------------------------------------------------

\begin{center}
\begin{tabularx}{\linewidth}{lXl}
\toprule
\textbf{Ratio} & \textbf{Formula} & \textbf{Measures} \\
\midrule
\multicolumn{3}{l}{\textit{Profitability}} \\
Gross Margin     & Gross Profit / Revenue                          & Pricing power \\
EBITDA Margin    & EBITDA / Revenue                                & Operating efficiency \\
ROE              & Net Income / Avg. Equity                        & Return to shareholders \\
ROA              & Net Income / Avg. Assets                        & Asset efficiency \\
\midrule
\multicolumn{3}{l}{\textit{Leverage}} \\
Debt/Equity      & Total Debt / Total Equity                       & Capital structure \\
Interest Coverage & EBIT / Interest Expense                        & Ability to service debt \\
\midrule
\multicolumn{3}{l}{\textit{Liquidity}} \\
Current Ratio    & Current Assets / Current Liabilities            & Short-term solvency \\
Quick Ratio      & (CA $-$ Inventory) / CL                        & Stringent liquidity \\
\midrule
\multicolumn{3}{l}{\textit{Efficiency}} \\
Asset Turnover   & Revenue / Avg. Assets                           & Asset productivity \\
Inventory Days   & (Inventory / COGS) $\times$ 365                 & Inventory management \\
\bottomrule
\end{tabularx}
\end{center}

\rmkb{\textbf{DuPont decomposition}: $ROE = \underbrace{\frac{\text{NI}}{\text{Sales}}}_{\text{margin}} \times \underbrace{\frac{\text{Sales}}{\text{Assets}}}_{\text{turnover}} \times \underbrace{\frac{\text{Assets}}{\text{Equity}}}_{\text{leverage}}$.}

% --------------------------------------------------------------------------
\section{Accrual Accounting vs.\ Cash Accounting}
% --------------------------------------------------------------------------

\defn{Accrual Basis}{
    Revenue is recognised when \textbf{earned} (goods delivered/services rendered),
    not when cash is received.  Expenses are recognised when \textbf{incurred},
    not when paid.  Result: net income $\neq$ cash flow in any given period.
    Key accruals:
    \begin{itemize}
        \item \textbf{Accounts Receivable}: revenue earned, cash not yet received.
        \item \textbf{Accounts Payable}: expense incurred, cash not yet paid.
        \item \textbf{Deferred Revenue}: cash received, revenue not yet earned.
        \item \textbf{Prepaid Expense}: cash paid, expense not yet incurred.
    \end{itemize}
}

% ==============================================================================
\part{Corporate Finance}
% ==============================================================================

% ------------------------------------------------------------------------------
\chapter{Time Value of Money}
% ------------------------------------------------------------------------------

\section{Discounting and Compounding}

\defn{Present and Future Value}{
    A cash flow $C$ received in $n$ periods, discounted at rate $r$ per period:
    \[
        PV = \frac{C}{(1+r)^n}, \qquad FV = C \cdot (1+r)^n.
    \]
    \textbf{Continuous compounding}: $PV = C\,e^{-rn}$, $FV = C\,e^{rn}$.
}

\defn{Annuities and Perpetuities}{
    \textbf{Ordinary annuity} (payments at end of period):
    \[
        PV_{\text{annuity}} = PMT \cdot \frac{1 - (1+r)^{-n}}{r}.
    \]
    \textbf{Growing annuity} (payment grows at rate $g$, $g \neq r$):
    \[
        PV = PMT \cdot \frac{1 - [(1+g)/(1+r)]^n}{r - g}.
    \]
    \textbf{Perpetuity}: $PV = PMT/r$.
    \textbf{Growing perpetuity} (Gordon Growth Model): $PV = PMT/(r-g)$, $g < r$.
}

\defn{Effective Annual Rate}{
    For $m$ compounding periods per year at nominal rate $r_{\text{nom}}$:
    \[
        EAR = \left(1 + \frac{r_{\text{nom}}}{m}\right)^m - 1.
    \]
    Continuous: $EAR = e^{r_{\text{nom}}} - 1$.
}

% ------------------------------------------------------------------------------
\chapter{Capital Budgeting}
% ------------------------------------------------------------------------------

\section{Investment Decision Rules}

\defn{Net Present Value}{
    \[
        NPV = \sum_{t=0}^{T} \frac{CF_t}{(1+r)^t} = -C_0 + \sum_{t=1}^{T} \frac{CF_t}{(1+r)^t}.
    \]
    \textbf{NPV Rule}: accept project iff $NPV > 0$.  NPV correctly accounts for
    the time value of money, risk, and opportunity cost of capital.
}

\defn{Internal Rate of Return}{
    The IRR solves $NPV(IRR) = 0$:
    \[
        0 = \sum_{t=0}^{T} \frac{CF_t}{(1+IRR)^t}.
    \]
    \textbf{IRR Rule}: accept iff $IRR > $ hurdle rate (cost of capital).
    Pitfalls: (1) multiple IRRs if sign changes more than once; (2) scale problem
    (higher IRR $\not\Rightarrow$ higher NPV); (3) assumes reinvestment at IRR.
    Always prefer NPV for mutually exclusive projects.
}

\rmkb{Other metrics: \textbf{Payback Period} (time to recoup $C_0$) ignores TVM
and post-payback flows.  \textbf{Profitability Index} $= NPV/C_0$ useful for
capital rationing (ranking projects under budget constraint).}

% ------------------------------------------------------------------------------
\chapter{Risk and Cost of Capital}
% ------------------------------------------------------------------------------

\section{Weighted Average Cost of Capital}

\defn{WACC}{
    The after-tax required return on the firm's assets, weighted by capital structure:
    \[
        WACC = \frac{E}{V} r_E + \frac{D}{V} r_D (1 - \tau_C),
    \]
    where $V = E + D$ is total firm value, $r_E$ is cost of equity, $r_D$ is
    pre-tax cost of debt, and $\tau_C$ is the corporate tax rate.
    Use WACC to discount FCFs when capital structure is stable.
}

\defn{Cost of Equity via CAPM}{
    \[
        r_E = r_f + \beta_E \cdot [E(R_M) - r_f].
    \]
    $\beta_E$ (equity beta) incorporates both business risk and financial leverage.
    \textbf{Unlevering}: $\beta_A = \frac{E}{E+D(1-\tau_C)} \beta_E + \frac{D(1-\tau_C)}{E+D(1-\tau_C)} \beta_D \approx \frac{E}{E+D}\beta_E$ (if $\beta_D \approx 0$).
}

% ------------------------------------------------------------------------------
\chapter{Capital Structure}
% ------------------------------------------------------------------------------

\section{Modigliani--Miller Theorems}

\thm{MM Proposition I (No Taxes)}{
    In perfect capital markets (no taxes, no bankruptcy costs, no information
    asymmetry, no transaction costs), the value of a firm is independent of its
    capital structure:
    \[
        V_L = V_U.
    \]
}

\thm{MM Proposition II (No Taxes)}{
    The cost of equity is an increasing linear function of leverage:
    \[
        r_E = r_A + \frac{D}{E}(r_A - r_D),
    \]
    where $r_A$ is the unlevered cost of equity.  WACC is constant regardless
    of leverage.
}

\thm{MM with Corporate Taxes}{
    Interest payments are tax-deductible, creating a \textbf{tax shield}:
    \[
        V_L = V_U + \tau_C D \quad \text{(assuming perpetual debt)}.
    \]
    Implication: firms should use maximum leverage.  Real-world limits:
    financial distress costs and agency costs.
}

\defn{Trade-Off Theory}{
    Optimal capital structure balances the tax shield benefit of debt against
    financial distress costs:
    \[
        V_L = V_U + PV(\text{Tax Shield}) - PV(\text{Distress Costs}).
    \]
    Predicts an interior optimal leverage ratio for each firm.
}

\defn{Pecking Order Theory}{
    Firms prefer internal financing (retained earnings), then debt, then equity.
    Rationale: adverse selection costs are lowest for internal funds, moderate
    for debt, and highest for equity (Myers-Majluf).
}

% ==============================================================================
\part{Investments and Portfolio Theory}
% ==============================================================================

% ------------------------------------------------------------------------------
\chapter{Financial Markets and Return Measurement}
% ------------------------------------------------------------------------------

\section{Asset Classes and Market Structure}

\defn{Primary vs.\ Secondary Markets}{
    \textbf{Primary market}: new securities are issued (IPO, seasoned offering).
    \textbf{Secondary market}: previously issued securities trade between investors
    (NYSE, NASDAQ).  Market efficiency implies secondary prices reflect all
    available information.
}

\section{Return Measurement}

\defn{Holding Period Return}{
    For a single period with beginning price $P_0$, ending price $P_1$,
    and dividend $D_1$:
    \[
        HPR = \frac{P_1 - P_0 + D_1}{P_0} = \frac{P_1 + D_1}{P_0} - 1.
    \]
}

\defn{Multi-Period Returns}{
    \textbf{Arithmetic mean} (best estimate of expected single-period return):
    \[
        \bar{R}_A = \frac{1}{T}\sum_{t=1}^T R_t.
    \]
    \textbf{Geometric mean} (best estimate of compound growth rate):
    \[
        \bar{R}_G = \left(\prod_{t=1}^T (1+R_t)\right)^{1/T} - 1.
    \]
    Always $\bar{R}_G \leq \bar{R}_A$; equality iff all returns are equal.
    Approximation: $\bar{R}_G \approx \bar{R}_A - \sigma^2/2$.
}

\defn{Risk Measures}{
    Sample variance and standard deviation of returns:
    \[
        s^2 = \frac{1}{T-1}\sum_{t=1}^T (R_t - \bar{R})^2, \qquad
        s = \sqrt{s^2}.
    \]
    \textbf{Sharpe Ratio}: $SR = (E[R] - r_f)/\sigma$ (reward per unit of total risk).
}

% ------------------------------------------------------------------------------
\chapter{Portfolio Theory}
% ------------------------------------------------------------------------------

\section{Mean-Variance Analysis}

\defn{Portfolio Return and Variance}{
    For $n$ assets with weights $\mathbf{w} = (w_1,\ldots,w_n)^\top$, $\mathbf{1}^\top\mathbf{w} = 1$:
    \[
        E[R_p] = \mathbf{w}^\top \boldsymbol{\mu} = \sum_{i=1}^n w_i \mu_i,
    \]
    \[
        \sigma_p^2 = \mathbf{w}^\top \boldsymbol{\Sigma}\, \mathbf{w}
        = \sum_{i=1}^n\sum_{j=1}^n w_i w_j \sigma_{ij},
    \]
    where $\boldsymbol{\Sigma}$ is the $n\times n$ covariance matrix and
    $\sigma_{ij} = \rho_{ij}\sigma_i\sigma_j$.
}

\defn{Diversification}{
    For an equally weighted portfolio of $n$ assets with average variance
    $\bar{\sigma}^2$ and average covariance $\bar{\sigma}_{ij}$:
    \[
        \sigma_p^2 = \frac{1}{n}\bar{\sigma}^2 + \frac{n-1}{n}\bar{\sigma}_{ij}
        \;\xrightarrow{n\to\infty}\; \bar{\sigma}_{ij}.
    \]
    Idiosyncratic (diversifiable) risk vanishes; systematic risk remains.
}

\section{Efficient Frontier}

\defn{Minimum Variance Portfolio}{
    The \textbf{Global Minimum Variance} (GMV) portfolio solves:
    \[
        \min_{\mathbf{w}}\; \mathbf{w}^\top \boldsymbol{\Sigma}\, \mathbf{w}
        \quad \text{s.t.} \quad \mathbf{1}^\top \mathbf{w} = 1.
    \]
    Solution: $\mathbf{w}_{GMV} = \frac{\boldsymbol{\Sigma}^{-1}\mathbf{1}}{\mathbf{1}^\top\boldsymbol{\Sigma}^{-1}\mathbf{1}}$.
    Adding the return target $\mathbf{w}^\top\boldsymbol{\mu} = \mu_p$ traces the
    \textbf{efficient frontier} in $(\sigma, E[R])$ space.
}

\defn{Capital Allocation Line and Tangency Portfolio}{
    With a risk-free asset at rate $r_f$, the investor's opportunity set becomes
    the \textbf{Capital Allocation Line (CAL)}:
    \[
        E[R_p] = r_f + \frac{E[R_T] - r_f}{\sigma_T}\,\sigma_p,
    \]
    where $T$ is the \textbf{tangency portfolio} maximising the Sharpe ratio:
    \[
        \mathbf{w}_T = \frac{\boldsymbol{\Sigma}^{-1}(\boldsymbol{\mu} - r_f\mathbf{1})}
        {\mathbf{1}^\top\boldsymbol{\Sigma}^{-1}(\boldsymbol{\mu} - r_f\mathbf{1})}.
    \]
}

\thm{Two-Fund Separation Theorem}{
    Every mean-variance efficient portfolio is a linear combination of the
    risk-free asset and the tangency portfolio.  All investors, regardless of
    risk aversion, choose a portfolio on the CAL through $T$.
}

\defn{Investor Utility and Optimal Allocation}{
    For mean-variance utility $U = E[R_p] - \frac{A}{2}\sigma_p^2$, the optimal
    fraction $w^*$ invested in the risky asset (or risky portfolio) satisfies:
    \[
        w^* = \frac{E[R] - r_f}{A\,\sigma^2}.
    \]
    Higher risk aversion ($A$) or higher variance $\Rightarrow$ lower allocation to risky assets.
}

% ------------------------------------------------------------------------------
\chapter{CAPM and Factor Models}
% ------------------------------------------------------------------------------

\section{Capital Asset Pricing Model}

\defn{CAPM}{
    In equilibrium, all investors hold the market portfolio $M$.  For any asset $i$:
    \[
        E[R_i] = r_f + \beta_i\,(E[R_M] - r_f),
    \]
    where the \textbf{beta} measures systematic risk:
    \[
        \beta_i = \frac{\text{Cov}(R_i, R_M)}{\text{Var}(R_M)}
        = \frac{\sigma_{iM}}{\sigma_M^2}.
    \]
    The \textbf{Security Market Line (SML)} plots $E[R_i]$ against $\beta_i$.
    The \textbf{Capital Market Line (CML)} plots $E[R_p]$ against $\sigma_p$ for
    efficient (fully diversified) portfolios:
    \[
        E[R_p] = r_f + \frac{E[R_M] - r_f}{\sigma_M}\,\sigma_p.
    \]
}

\defn{Risk Decomposition and Alpha}{
    Total variance decomposes into systematic and idiosyncratic parts:
    \[
        \sigma_i^2 = \beta_i^2\sigma_M^2 + \sigma_{\varepsilon_i}^2.
    \]
    \textbf{Alpha} (Jensen's alpha) measures risk-adjusted outperformance:
    \[
        \alpha_i = E[R_i] - \left[r_f + \beta_i(E[R_M] - r_f)\right].
    \]
    In an efficient market, $\alpha = 0$ for all assets.
}

\section{Factor Models}

\defn{Fama-French Three-Factor Model}{
    Extends CAPM with size and value factors:
    \[
        E[R_i] - r_f = \alpha_i + \beta_{iM}(R_M - r_f) + \beta_{iSMB}\cdot SMB + \beta_{iHML}\cdot HML,
    \]
    where $SMB$ (Small Minus Big) captures the size premium, and $HML$ (High Minus Low
    book-to-market) captures the value premium.  Carhart adds $MOM$ (momentum).
}

\rmkb{The CAPM assumptions that fail empirically: (1) all investors hold the same
mean-variance efficient portfolio; (2) no taxes or transaction costs; (3) unlimited
borrowing at $r_f$.  Factor models are more descriptive; they do not necessarily
imply risk-based explanations (may reflect mispricing or data mining).}

% ------------------------------------------------------------------------------
\chapter{Equity Valuation}
% ------------------------------------------------------------------------------

\section{Dividend Discount Model}

\defn{DDM (Gordon Growth Model)}{
    The price of a stock equals the PV of all future dividends:
    \[
        P_0 = \sum_{t=1}^{\infty} \frac{D_t}{(1+r_E)^t}.
    \]
    For constant growth $g$ (Gordon Growth Model):
    \[
        P_0 = \frac{D_1}{r_E - g} = \frac{D_0(1+g)}{r_E - g}, \qquad g < r_E.
    \]
    Implies: $r_E = D_1/P_0 + g$ (dividend yield + capital gains yield).
}

\section{Free Cash Flow to Firm (FCFF) Model}

\defn{DCF Valuation}{
    Enterprise value equals the PV of future FCFs discounted at WACC:
    \[
        EV = \sum_{t=1}^{T} \frac{FCFF_t}{(1+WACC)^t} + \frac{TV_T}{(1+WACC)^T},
    \]
    where terminal value $TV_T = FCFF_{T+1}/(WACC - g)$ (Gordon Growth).
    \[
        \text{Equity Value} = EV - \text{Net Debt}.
    \]
    $FCFF = EBIT(1-\tau) + D\&A - \Delta NWC - \text{Capex}$.
}

\section{Relative Valuation}

\begin{center}
\begin{tabularx}{\linewidth}{lXX}
\toprule
\textbf{Multiple} & \textbf{Formula} & \textbf{Use} \\
\midrule
P/E       & Price per share / EPS         & Mature, profitable firms \\
EV/EBITDA & Enterprise Value / EBITDA     & Cross-capital-structure comparisons \\
EV/Sales  & Enterprise Value / Revenue    & High-growth or loss-making firms \\
P/B       & Price / Book Value per share  & Financials; asset-heavy sectors \\
PEG       & P/E / Earnings Growth Rate    & Growth-adjusted P/E \\
\bottomrule
\end{tabularx}
\end{center}

\rmkb{Multiples are only meaningful relative to peers.  Always adjust for differences
in growth, risk, and margins across comparables.}

% ------------------------------------------------------------------------------
\chapter{Fixed Income}
% ------------------------------------------------------------------------------

\section{Bond Pricing}

\defn{Bond Price}{
    A bond with face value $F$, coupon rate $c$, maturity $T$, and yield to maturity
    $y$ has price:
    \[
        P = \sum_{t=1}^{T} \frac{cF}{(1+y)^t} + \frac{F}{(1+y)^T}
        = cF\cdot\frac{1-(1+y)^{-T}}{y} + F(1+y)^{-T}.
    \]
    $P > F$ iff $y < c$ (premium bond); $P < F$ iff $y > c$ (discount bond).
    Price and yield move \textbf{inversely}.
}

\section{Duration and Convexity}

\defn{Macaulay Duration}{
    The weighted-average time to receipt of cash flows (weights = PV of each
    cash flow as fraction of total price):
    \[
        D_{Mac} = \frac{1}{P}\sum_{t=1}^{T} t \cdot \frac{CF_t}{(1+y)^t}.
    \]
    \textbf{Modified Duration}: $D_{Mod} = D_{Mac}/(1+y)$.
    Price sensitivity: $\Delta P \approx -D_{Mod} \cdot P \cdot \Delta y$.
}

\defn{Convexity}{
    Second-order correction to the price-yield relationship:
    \[
        \frac{\Delta P}{P} \approx -D_{Mod}\,\Delta y + \frac{1}{2}\,C\,(\Delta y)^2,
    \]
    where convexity $C = \frac{1}{P}\frac{d^2P}{dy^2}$.
    Higher convexity is always beneficial: the bond rises more (falls less) than
    the duration approximation predicts for any yield change.
}

\section{Yield Curve}

\defn{Term Structure of Interest Rates}{
    The \textbf{yield curve} plots spot rates $s_t$ (YTM of zero-coupon bonds)
    against maturity $t$.  Shapes: normal (upward sloping), flat, inverted.
    \textbf{Forward rate} $f_{t,t+1}$ is the implied one-period rate from $t$ to $t+1$:
    \[
        (1+s_{t+1})^{t+1} = (1+s_t)^t\,(1+f_{t,t+1}).
    \]
    Theories: (1) Expectations: $f_{t,t+1} = E[s_{t+1}]$;
    (2) Liquidity preference: $f_{t,t+1} > E[s_{t+1}]$ (term premium);
    (3) Market segmentation: supply/demand at each maturity.
}

% ------------------------------------------------------------------------------
\chapter{Derivatives and Options}
% ------------------------------------------------------------------------------

\section{Option Payoffs and Strategies}

\defn{Call and Put Options}{
    A \textbf{European call} gives the right to buy asset $S$ at strike $K$ at
    maturity $T$. Payoff at expiry: $\max(S_T - K, 0)$.
    A \textbf{European put}: payoff $\max(K - S_T, 0)$.
    Premium (price) paid upfront; net profit = payoff $-$ premium.
}

\defn{Put--Call Parity}{
    For European options on a non-dividend-paying stock:
    \[
        C - P = S_0 - K e^{-rT},
    \]
    where $C$ and $P$ are call and put prices with the same $K$ and $T$.
    Violation implies arbitrage.
}

\begin{center}
\begin{tabularx}{\linewidth}{lXl}
\toprule
\textbf{Strategy} & \textbf{Construction} & \textbf{View} \\
\midrule
Covered Call     & Long stock $+$ Short call           & Mildly bullish, income \\
Protective Put   & Long stock $+$ Long put             & Bullish, insured \\
Bull Spread      & Long call($K_1$) $+$ Short call($K_2$), $K_1 < K_2$ & Moderately bullish \\
Bear Spread      & Long put($K_2$) $+$ Short put($K_1$) & Moderately bearish \\
Straddle         & Long call $+$ Long put (same $K,T$) & High volatility \\
Strangle         & Long call($K_H$) $+$ Long put($K_L$), $K_H > K_L$ & High vol, cheaper \\
\bottomrule
\end{tabularx}
\end{center}

\section{Black-Scholes-Merton Model}

\defn{BSM Price}{
    Assuming the stock price follows geometric Brownian motion
    $dS = \mu S\,dt + \sigma S\,dW_t$, the no-arbitrage price of a European call is:
    \[
        C = S_0 N(d_1) - K e^{-rT} N(d_2),
    \]
    \[
        d_1 = \frac{\ln(S_0/K) + (r + \sigma^2/2)T}{\sigma\sqrt{T}}, \qquad
        d_2 = d_1 - \sigma\sqrt{T},
    \]
    where $N(\cdot)$ is the standard normal CDF.
    Put price: $P = Ke^{-rT}N(-d_2) - S_0 N(-d_1)$.
}

\defn{The Greeks}{
    \begin{center}
    \begin{tabularx}{\linewidth}{llX}
    \toprule
    \textbf{Greek} & \textbf{Definition} & \textbf{Interpretation} \\
    \midrule
    Delta $\Delta$ & $\partial C/\partial S$        & Change in option price per \$1 change in $S$; $\in(0,1)$ for calls \\
    Gamma $\Gamma$ & $\partial^2 C/\partial S^2$    & Rate of change of delta; highest near ATM \\
    Theta $\Theta$ & $\partial C/\partial t$        & Time decay; typically negative for long options \\
    Vega $\nu$     & $\partial C/\partial \sigma$   & Sensitivity to volatility; positive for long options \\
    Rho $\rho$     & $\partial C/\partial r$        & Sensitivity to interest rate \\
    \bottomrule
    \end{tabularx}
    \end{center}
}

% ==============================================================================
\part{Quantitative Finance}
% ==============================================================================

% ------------------------------------------------------------------------------
\chapter{Portfolio Optimisation}
% ------------------------------------------------------------------------------

\section{Mean-Variance Optimisation}

\defn{Markowitz Optimisation Problem}{
    For a target return $\mu_p^*$, the minimum variance portfolio solves:
    \[
        \min_{\mathbf{w}}\; \mathbf{w}^\top \boldsymbol{\Sigma}\, \mathbf{w}
        \quad \text{s.t.} \quad
        \mathbf{w}^\top\boldsymbol{\mu} = \mu_p^*, \quad
        \mathbf{1}^\top\mathbf{w} = 1.
    \]
    Using Lagrangians $\lambda$ and $\gamma$, the KKT conditions give the closed form:
    \[
        \mathbf{w}^* = \lambda^* \boldsymbol{\Sigma}^{-1}\boldsymbol{\mu}
        + \gamma^* \boldsymbol{\Sigma}^{-1}\mathbf{1},
    \]
    where the multipliers solve the $2\times 2$ system:
    \[
        \begin{pmatrix} A & B \\ B & C \end{pmatrix}
        \begin{pmatrix} \lambda^* \\ \gamma^* \end{pmatrix}
        = \begin{pmatrix} \mu_p^* \\ 1 \end{pmatrix},
    \]
    with $A = \boldsymbol{\mu}^\top\boldsymbol{\Sigma}^{-1}\boldsymbol{\mu}$,
    $B = \boldsymbol{\mu}^\top\boldsymbol{\Sigma}^{-1}\mathbf{1}$,
    $C = \mathbf{1}^\top\boldsymbol{\Sigma}^{-1}\mathbf{1}$, and
    $\Delta = AC - B^2 > 0$.
    The minimum variance of the frontier portfolio at $\mu_p^*$ is:
    \[
        \sigma^{*2}(\mu_p^*) = \frac{C\mu_p^{*2} - 2B\mu_p^* + A}{\Delta}.
    \]
}

\section{Sharpe Ratio Maximisation}

\defn{Tangency Portfolio (Analytical Form)}{
    With a risk-free rate $r_f$, the excess return vector is
    $\tilde{\boldsymbol{\mu}} = \boldsymbol{\mu} - r_f\mathbf{1}$.
    The tangency (maximum Sharpe) portfolio is:
    \[
        \mathbf{w}_T = \frac{\boldsymbol{\Sigma}^{-1}\tilde{\boldsymbol{\mu}}}
        {\mathbf{1}^\top\boldsymbol{\Sigma}^{-1}\tilde{\boldsymbol{\mu}}},
        \qquad
        SR_T = \sqrt{\tilde{\boldsymbol{\mu}}^\top\boldsymbol{\Sigma}^{-1}\tilde{\boldsymbol{\mu}}}.
    \]
}

\section{Practical Constraints and Shrinkage}

\rmkb{The classical Markowitz solution is notoriously unstable: small errors in
$\boldsymbol{\mu}$ and $\boldsymbol{\Sigma}$ lead to extreme allocations.  Practical remedies:
\begin{itemize}
    \item \textbf{Long-only constraint}: $w_i \geq 0$ (convex QP, solved via active-set or interior-point methods).
    \item \textbf{Shrinkage estimators}: Ledoit-Wolf shrinks $\hat{\boldsymbol{\Sigma}}$ toward a structured target (e.g., diagonal) to reduce estimation error.
    \item \textbf{Black-Litterman}: blends market-implied equilibrium returns with investor views via Bayesian updating, producing better-conditioned portfolios.
    \item \textbf{Risk parity}: allocate equal risk contribution $w_i(\boldsymbol{\Sigma}\mathbf{w})_i / \sigma_p^2 = 1/n$ across assets.
\end{itemize}}

% ------------------------------------------------------------------------------
\chapter{Mathematical Finance}
% ------------------------------------------------------------------------------

\section{Single-Period Model and No-Arbitrage}

\defn{Single-Period Market Model}{
    A market consists of:
    \begin{itemize}
        \item A \textbf{risk-free bond}: value 1 at time 0, value $1+r$ at time 1.
        \item $d$ \textbf{risky assets}: price vector $S_0 \in \mathbb{R}^d$ at time 0;
              random payoff $S_1(\omega)$ at time 1 in state $\omega \in \Omega$
              (finite state space, $|\Omega| = k$).
        \item A \textbf{portfolio} $H = (H^0, H^1, \ldots, H^d)$: holdings of bond
              and risky assets.
    \end{itemize}
    Portfolio value: $V_0 = H^0 + H \cdot S_0$, $V_1(\omega) = H^0(1+r) + H \cdot S_1(\omega)$.
}

\defn{Arbitrage}{
    A portfolio $H$ is an \textbf{arbitrage} if:
    \[
        V_0 \leq 0, \quad V_1(\omega) \geq 0\ \forall\omega, \quad
        \exists\,\omega^*: V_1(\omega^*) > 0.
    \]
    (Get something from nothing, or nothing for free with positive probability of gain.)
}

\thm{Fundamental Theorem of Asset Pricing (FTAP) --- Single Period}{
    The market is \textbf{arbitrage-free} if and only if there exists a
    \textbf{risk-neutral (equivalent martingale) measure} $\mathbb{Q}$,
    i.e., a probability measure with $\mathbb{Q}(\omega) > 0\ \forall\omega$ such that
    \[
        S_0^i = \frac{1}{1+r}\,\mathbb{E}^{\mathbb{Q}}[S_1^i]
        \qquad \forall\, i = 1, \ldots, d.
    \]
    The market is \textbf{complete} (every payoff is replicable) iff $\mathbb{Q}$
    is unique.
}

\section{Binomial Model}

\defn{One-Period Binomial Model}{
    Stock price $S_0$ moves to $S_0 u$ (up) with physical probability $p$,
    or $S_0 d$ (down) with $1-p$, where $d < 1+r < u$ (no-arbitrage condition).
    The unique risk-neutral probability:
    \[
        q = \frac{(1+r) - d}{u - d}.
    \]
    Price of a derivative with payoffs $(V_u, V_d)$:
    \[
        V_0 = \frac{qV_u + (1-q)V_d}{1+r}.
    \]
    \textbf{Replicating portfolio}: $\Delta = (V_u - V_d)/(S_0 u - S_0 d)$ shares,
    bond position $B = (V_0 - \Delta S_0)/(1+r)^0$.
}

\defn{Multi-Period Binomial Tree}{
    After $n$ steps, each path is a sequence of ups and downs.
    The $T$-period European option price:
    \[
        V_0 = \frac{1}{(1+r)^T}\sum_{k=0}^{T}\binom{T}{k}q^k(1-q)^{T-k} h(S_0 u^k d^{T-k}),
    \]
    where $h(S_T)$ is the payoff function.
    American options: use backward induction, taking $\max(\text{intrinsic value}, \text{continuation value})$ at each node.
}

% ------------------------------------------------------------------------------
\chapter{Stochastic Processes and Black-Scholes}
% ------------------------------------------------------------------------------

\section{Brownian Motion}

\defn{Standard Brownian Motion}{
    A stochastic process $\{W_t\}_{t \geq 0}$ is a \textbf{standard Brownian motion}
    (Wiener process) if:
    \begin{enumerate}
        \item $W_0 = 0$ a.s.
        \item \textbf{Independent increments}: for $0 \leq s < t$, $W_t - W_s$ is
              independent of $\mathcal{F}_s$ (the natural filtration).
        \item \textbf{Gaussian increments}: $W_t - W_s \sim \mathcal{N}(0, t-s)$.
        \item \textbf{Continuous paths}: $t \mapsto W_t$ is continuous a.s.
    \end{enumerate}
}

\rmkb{Key properties of Brownian motion:
\begin{itemize}
    \item $E[W_t] = 0$, $\text{Var}(W_t) = t$, $\text{Cov}(W_s, W_t) = \min(s,t)$.
    \item Paths are continuous but nowhere differentiable (a.s.).
    \item \textbf{Quadratic variation}: $[W]_t = t$ (i.e., $\sum (W_{t_i} - W_{t_{i-1}})^2 \to t$).
    \item $W_t$ is a martingale: $E[W_t | \mathcal{F}_s] = W_s$.
\end{itemize}}

\section{It\^{o} Calculus}

\defn{It\^{o} Integral}{
    For an adapted process $\{H_t\}$, the \textbf{It\^{o} integral} is defined as:
    \[
        \int_0^T H_t\,dW_t = \lim_{n\to\infty} \sum_{i=0}^{n-1} H_{t_i}(W_{t_{i+1}} - W_{t_i}),
    \]
    using the \emph{left-endpoint} (non-anticipating) evaluation.  This differs
    from the Stratonovich integral and preserves the martingale property.
}

\thm{It\^{o}'s Lemma}{
    Let $X_t$ be an It\^{o} process: $dX_t = \mu_t\,dt + \sigma_t\,dW_t$.
    For $f \in C^{2,1}(\mathbb{R}\times[0,T])$:
    \[
        df(X_t, t) = \left(\frac{\partial f}{\partial t} + \mu_t\frac{\partial f}{\partial x}
        + \frac{\sigma_t^2}{2}\frac{\partial^2 f}{\partial x^2}\right)dt
        + \sigma_t\frac{\partial f}{\partial x}\,dW_t.
    \]
    The extra term $\frac{\sigma^2}{2}\frac{\partial^2 f}{\partial x^2}dt$ arises from
    the quadratic variation of $W_t$: $(dW_t)^2 = dt$ (It\^{o}'s table:
    $dt \cdot dt = 0$, $dt \cdot dW = 0$, $dW \cdot dW = dt$).
}

\section{Geometric Brownian Motion and the BSM Equation}

\defn{Geometric Brownian Motion}{
    The standard model for stock prices:
    \[
        dS_t = \mu S_t\,dt + \sigma S_t\,dW_t.
    \]
    By It\^{o}'s lemma applied to $f(S_t) = \ln S_t$:
    \[
        d\ln S_t = \left(\mu - \frac{\sigma^2}{2}\right)dt + \sigma\,dW_t,
    \]
    so $\ln S_t \sim \mathcal{N}\!\left(\ln S_0 + (\mu-\sigma^2/2)t,\; \sigma^2 t\right)$
    and:
    \[
        S_t = S_0 \exp\!\left[\left(\mu - \frac{\sigma^2}{2}\right)t + \sigma W_t\right].
    \]
}

\thm{Black-Scholes PDE}{
    Consider a self-financing portfolio in stock $S$ and bond $B = e^{rt}$
    that replicates derivative payoff $V(S, t)$.  The no-arbitrage condition yields
    the \textbf{Black-Scholes PDE}:
    \[
        \frac{\partial V}{\partial t} + rS\frac{\partial V}{\partial S}
        + \frac{1}{2}\sigma^2 S^2\frac{\partial^2 V}{\partial S^2} - rV = 0,
    \]
    with terminal condition $V(S,T) = h(S)$ (payoff function).
}

\pf{
    (Delta-hedging derivation.)  Let $\Pi = V - \Delta S$ be a hedged portfolio.
    By It\^{o}'s lemma:
    \[
        d\Pi = \left(\frac{\partial V}{\partial t} + \frac{1}{2}\sigma^2 S^2
        \frac{\partial^2 V}{\partial S^2}\right)dt
        + \left(\frac{\partial V}{\partial S} - \Delta\right)\!dS.
    \]
    Setting $\Delta = \partial V/\partial S$ eliminates the stochastic term.
    The resulting riskless portfolio must earn the risk-free rate: $d\Pi = r\Pi\,dt$,
    giving the BS PDE.
}

\defn{Risk-Neutral Pricing and the Girsanov Theorem}{
    Under the physical measure $\mathbb{P}$, $dS = \mu S\,dt + \sigma S\,dW_t$.
    Define the market price of risk $\theta = (\mu - r)/\sigma$.
    By \textbf{Girsanov's theorem}, changing measure to $\mathbb{Q}$ via
    Radon-Nikod\'ym derivative
    \[
        \frac{d\mathbb{Q}}{d\mathbb{P}} = \exp\!\left(-\theta W_T - \frac{\theta^2}{2}T\right),
    \]
    the process $\tilde{W}_t = W_t + \theta t$ is a $\mathbb{Q}$-Brownian motion and:
    \[
        dS = r S\,dt + \sigma S\,d\tilde{W}_t.
    \]
    Under $\mathbb{Q}$, the discounted price $e^{-rt}S_t$ is a martingale, and any
    derivative price is:
    \[
        V_0 = e^{-rT}\,\mathbb{E}^{\mathbb{Q}}[h(S_T)].
    \]
    For a European call, this integral yields the BSM formula.
}

\thm{BSM Formula (Derivation Sketch)}{
    Under $\mathbb{Q}$: $S_T = S_0\exp[(r - \sigma^2/2)T + \sigma\tilde{W}_T]$,
    $\tilde{W}_T \sim \mathcal{N}(0,T)$.
    \begin{align*}
        C &= e^{-rT}\mathbb{E}^{\mathbb{Q}}[\max(S_T - K, 0)] \\
          &= e^{-rT}\mathbb{E}^{\mathbb{Q}}[(S_T - K)\mathbf{1}_{S_T > K}] \\
          &= S_0\,\mathbb{Q}(S_T > K \mid S_0)
             - Ke^{-rT}\,\mathbb{Q}(S_T > K \mid S_0) \\
          &= S_0 N(d_1) - Ke^{-rT}N(d_2),
    \end{align*}
    where the two probabilities $N(d_1)$ and $N(d_2)$ are computed by standardising
    the lognormal $S_T$ under $\mathbb{Q}$ and the $S$-forward measure respectively.
}

% ==============================================================================
\part{Game Theory}
% ==============================================================================

% ------------------------------------------------------------------------------
\chapter{Normal Form Games}
% ------------------------------------------------------------------------------

\section{Strategic Form and Dominance}

\defn{Normal Form Game}{
    A game $G = (N, \{S_i\}_{i\in N}, \{u_i\}_{i\in N})$ consists of:
    \begin{itemize}
        \item A finite set of \textbf{players} $N = \{1, \ldots, n\}$.
        \item For each player $i$: a \textbf{strategy set} $S_i$
              (finite for matrix games).
        \item \textbf{Payoff functions} $u_i: S_1 \times \cdots \times S_n \to \mathbb{R}$.
    \end{itemize}
    The strategy profile $s = (s_1, \ldots, s_n) \in S = \prod_i S_i$.
    Notation: $s_{-i}$ denotes all players' strategies except $i$.
}

\defn{Dominant Strategy}{
    Strategy $s_i^*$ \textbf{strictly dominates} $s_i'$ if
    $u_i(s_i^*, s_{-i}) > u_i(s_i', s_{-i})$ for all $s_{-i} \in S_{-i}$.
    \textbf{Iterated Elimination of Strictly Dominated Strategies (IESDS)}:
    repeatedly remove dominated strategies; the order does not matter;
    any rationalizable strategy survives IESDS.
}

\defn{Nash Equilibrium}{
    A strategy profile $s^* \in S$ is a \textbf{Nash Equilibrium (NE)} if for all
    players $i$ and all $s_i \in S_i$:
    \[
        u_i(s_i^*, s_{-i}^*) \geq u_i(s_i, s_{-i}^*).
    \]
    No player can profitably deviate given the others' strategies.
}

\defn{Mixed Strategy Nash Equilibrium}{
    A \textbf{mixed strategy} $\sigma_i \in \Delta(S_i)$ is a probability
    distribution over $S_i$.  Player $i$'s expected payoff:
    $u_i(\sigma) = \sum_{s\in S} \sigma_1(s_1)\cdots\sigma_n(s_n)\,u_i(s)$.
    In a \textbf{mixed strategy NE}, each player is indifferent among all
    strategies in the support of $\sigma_i^*$:
    \[
        u_i(s_i, \sigma_{-i}^*) = \text{const} \quad \forall s_i \in \text{supp}(\sigma_i^*).
    \]
}

\thm{Nash Existence Theorem}{
    Every finite game (finite $N$, finite $S_i$) has at least one Nash equilibrium
    in mixed strategies.
}

\pf{(Sketch.) Define the best-response correspondence $BR_i(\sigma_{-i}) = \arg\max_{\sigma_i} u_i(\sigma_i, \sigma_{-i})$.  The joint $BR(\sigma) = \prod_i BR_i(\sigma_{-i})$ maps the compact convex set $\prod_i \Delta(S_i)$ to itself.  Kakutani's fixed-point theorem applies (BR is upper hemicontinuous with non-empty convex values), guaranteeing a fixed point $\sigma^* = BR(\sigma^*)$, which is a NE.}

\section{Classic Examples}

\ex{
\textbf{Prisoner's Dilemma}: both players choose $\{C, D\}$; payoff matrix:
\[
\begin{array}{c|cc}
 & C & D \\\hline
C & (2,2) & (0,3) \\
D & (3,0) & (1,1)
\end{array}
\]
$D$ strictly dominates $C$ for both players.  Unique NE: $(D,D)$ with payoffs
$(1,1)$ --- Pareto dominated by $(C,C)$.  Illustrates the tension between
individual rationality and social efficiency.
}

\ex{
\textbf{Battle of the Sexes}: two pure NE (Opera, Opera) and (Football, Football),
plus one mixed NE where the mixing probabilities balance payoffs.
}

% ------------------------------------------------------------------------------
\chapter{Extensive Form Games}
% ------------------------------------------------------------------------------

\section{Sequential Games and Backward Induction}

\defn{Extensive Form Game}{
    Represented as a game tree with:
    \begin{itemize}
        \item \textbf{Nodes}: decision nodes (assigned to a player) and terminal nodes.
        \item \textbf{Information sets}: sets of nodes a player cannot distinguish.
        \item \textbf{Actions}: edges from a decision node.
        \item \textbf{Payoffs}: at each terminal node.
    \end{itemize}
    A \textbf{strategy} for player $i$ specifies an action at every information set.
}

\defn{Subgame Perfect Equilibrium (SPE)}{
    A strategy profile is an SPE if it induces a Nash equilibrium in every
    \textbf{subgame} (subtree starting from a single node with complete information).
    SPE eliminates non-credible threats.
    \textbf{Backward induction}: solve the game tree from terminal nodes upward,
    choosing the optimal action at each node --- yields the unique SPE in finite
    games of perfect information.
}

\ex{
\textbf{Stackelberg Duopoly}: leader (L) moves first, follower (F) observes and
responds.  Backward induction: F's best response $q_F^{BR}(q_L)$ is computed first;
L maximises profit anticipating F's response.  Result: L produces more than Cournot;
F produces less.  First-mover advantage arises from commitment.
}

% ------------------------------------------------------------------------------
\chapter{Repeated Games}
% ------------------------------------------------------------------------------

\section{Finitely and Infinitely Repeated Games}

\defn{Repeated Game}{
    The stage game $G$ is played at $t = 1, 2, \ldots$ Players observe all
    past actions and discount future payoffs by $\delta \in (0,1)$ per period.
    Total payoff: $\sum_{t=0}^{\infty}\delta^t u_i(a^t)$.
}

\thm{Backward Induction in Finite Repetition}{
    If the stage game has a unique NE, then in the finitely repeated game the unique
    SPE is to play the stage NE at every period.  Cooperation cannot be sustained
    by the threat of future punishment (unravelling from the last period).
}

\thm{Folk Theorem (Informal)}{
    In an infinitely repeated game, if players are sufficiently patient (high $\delta$),
    any feasible individually rational payoff vector can be supported as a NE
    payoff (and under stronger conditions, as an SPE payoff via \textbf{grim trigger}
    or \textbf{tit-for-tat} strategies).
}

\pf{(Grim Trigger example, Prisoner's Dilemma.)
Cooperate while the opponent has always cooperated; defect forever otherwise.
Cooperation is sustainable if the one-shot gain from defecting is outweighed
by the loss of future cooperation:
\[
    3 + \frac{\delta}{1-\delta}\cdot 1 \leq 2 + \frac{\delta}{1-\delta}\cdot 2
    \quad\Longleftrightarrow\quad \delta \geq \frac{1}{2}.
\]
For $\delta \geq 1/2$, cooperation is a subgame-perfect equilibrium.}

% ------------------------------------------------------------------------------
\chapter{Bayesian Games and Auctions}
% ------------------------------------------------------------------------------

\section{Games with Incomplete Information}

\defn{Bayesian Game}{
    A game where each player has a \textbf{type} $\theta_i \in \Theta_i$ drawn from
    a joint distribution $F(\theta)$, privately observed.  Player $i$ maximises
    expected payoff over others' types:
    \[
        \max_{a_i(\theta_i)}\; \mathbb{E}_{\theta_{-i}}[u_i(a_i, a_{-i}(\theta_{-i}), \theta)].
    \]
    A \textbf{Bayesian Nash Equilibrium (BNE)} is a strategy profile
    $a_i^*(\cdot): \Theta_i \to A_i$ such that each player's strategy is a best
    response to the others' strategies, given beliefs about others' types.
}

\section{Auction Theory}

\defn{Auction Formats}{
    \begin{center}
    \begin{tabularx}{\linewidth}{lX}
    \toprule
    \textbf{Format} & \textbf{Mechanism} \\
    \midrule
    English (ascending) & Open, price rises until one bidder remains; winner pays final price \\
    Dutch (descending)  & Price falls until first bidder accepts; strategically equivalent to FPSB \\
    First-Price Sealed Bid (FPSB) & Bidders submit sealed bids; highest bid wins, pays own bid \\
    Second-Price Sealed Bid (Vickrey) & Highest bid wins, pays second-highest bid \\
    \bottomrule
    \end{tabularx}
    \end{center}
}

\thm{Revenue Equivalence Theorem}{
    Suppose bidders have independent private values drawn from a common distribution
    $F$ on $[0, \bar{v}]$, are risk-neutral, and the auction allocates the object
    to the highest-value bidder with a zero expected payment for the lowest type.
    Then \textbf{all four standard auction formats yield the same expected revenue
    to the seller}.
}

\defn{FPSB Equilibrium Bid}{
    With $n$ bidders and values $v_i \sim U[0,\bar{v}]$ i.i.d., the symmetric
    Bayesian Nash equilibrium bid is:
    \[
        b^*(v) = \frac{n-1}{n}\,v.
    \]
    Each bidder shades their bid below their value; the shading diminishes as $n$ grows.
}

\defn{Vickrey Auction (SPSB)}{
    In the second-price auction, bidding own value $b_i(v_i) = v_i$ is a
    \textbf{weakly dominant strategy} for every bidder, regardless of the number of
    bidders or others' strategies.  Truth-telling is optimal because:
    \[
        \text{winning at second price} = \max_{j \neq i} b_j \leq v_i
        \Rightarrow \text{surplus} = v_i - \max_{j \neq i} b_j > 0.
    \]
    Overbidding risks winning at a price exceeding value; underbidding risks
    losing a profitable auction.
}

\rmkb{The Myerson optimal auction maximises expected revenue by introducing a
\textbf{reserve price} $r^*$: the seller withholds the good if the highest bid
falls below $r^*$.  With $n$ bidders and $v_i \sim U[0,1]$:
$r^* = 1/2$ regardless of $n$.}

\end{document}
